\documentclass[11pt]{article}
\usepackage[a4paper,margin=1in]{geometry}
\usepackage{amsmath,amssymb}
\usepackage{booktabs}
\usepackage{hyperref}

\title{Research Note: Worst-Case Guarantee of Demand-Weighted Greedy Coverage}
\author{}
\date{\today}

\begin{document}
\maketitle

\section{Problem Mapping}
Your demand-based ambulance objective is a \emph{weighted maximum coverage} objective.
Let each candidate ambulance location be a set, and each postal code be an element with nonnegative demand weight.
For a chosen set of locations $S$ with $|S| \le k$, the objective is
\[
F(S) = \sum_{e \in \cup_{A \in S} A} w_e,
\]
where $w_e \ge 0$ is demand.

This is a monotone submodular maximization problem under a cardinality constraint.

\section{Known Approximation Guarantee}
For monotone submodular maximization under a cardinality constraint, classic greedy gives
\[
F(S_{\text{greedy}}) \ge \left(1 - \left(1 - \frac{1}{k}\right)^k\right) F(S^*).
\]
Hence also
\[
F(S_{\text{greedy}}) \ge (1 - 1/e) F(S^*) \approx 0.632\, F(S^*).
\]

For small fixed $k$, the finite-$k$ bound is stronger:
\begin{itemize}
\item $k=2$: $1-(1-1/2)^2 = 3/4 = 0.75$.
\item $k=3$: $1-(2/3)^3 = 19/27 \approx 0.7037$.
\end{itemize}

\section{What This Means for Your Algorithm}
Your demand-greedy step picks the location with maximum \emph{marginal uncovered demand} each round.
So it is exactly the weighted maximum-coverage greedy rule.

The travel-time tie-break you added is used only when marginal demand is equal.
Therefore it does not weaken the approximation guarantee above.

\section{About Worst Cases}
Worst-case instances are known where greedy gets close to the bound and cannot be improved in polynomial time (unless standard complexity assumptions fail).
So seeing ratios much better than $1-1/e$ in random tests is normal; the bound is a worst-case lower bound, not an average-case prediction.

\section{Computational Search in Notebook}
An extra notebook cell was added to search for bad instances by:
\begin{enumerate}
\item Generating many random weighted coverage instances.
\item Computing greedy demand score.
\item Computing exact optimal score by brute force.
\item Reporting the worst observed ratio and the corresponding instance summary.
\end{enumerate}

This does not prove the tight theoretical worst case, but it helps find practical failure patterns for your implementation and data style.

\section{References}
\begin{enumerate}
\item G. L. Nemhauser, L. A. Wolsey, and M. L. Fisher, ``An analysis of approximations for maximizing submodular set functions I,'' \emph{Mathematical Programming}, 14, 265--294, 1978.
\item Maximum coverage problem (weighted version and greedy ratio), Wikipedia: \url{https://en.wikipedia.org/wiki/Maximum_coverage_problem}.
\item Submodular maximization under cardinality constraints, Wikipedia: \url{https://en.wikipedia.org/wiki/Submodular_set_function}.
\end{enumerate}

\end{document}
